\documentclass[12pt,a4paper]{article}

\usepackage[a4paper,margin=1in]{geometry}
\usepackage[T1]{fontenc}
\usepackage[utf8]{inputenc}
\usepackage{lmodern}
\usepackage{setspace}
\usepackage{graphicx}
\usepackage{enumitem}
\usepackage{array}
\usepackage{longtable}
\usepackage{booktabs}
\usepackage{tabularx}
\usepackage{titlesec}
\usepackage[hidelinks]{hyperref}
\usepackage{xcolor}
\usepackage{listings}

\setstretch{1.12}
\setlength{\parindent}{0pt}
\setlength{\parskip}{6pt}
\titleformat{\section}{\normalfont\large\bfseries}{\thesection}{1em}{}
\titleformat{\subsection}{\normalfont\normalsize\bfseries}{\thesubsection}{1em}{}

\lstset{
  language=SQL,
  basicstyle=\ttfamily\scriptsize,
  frame=single,
  breaklines=true,
  showstringspaces=false,
  columns=fullflexible
}

\newcolumntype{Y}{>{\raggedright\arraybackslash}X}

\begin{document}

\begin{center}
\textbf{UNIVERSITY OF CHITTAGONG}\\
Department of Computer Science and Engineering
\end{center}

\vspace{0.8cm}

\begin{center}
\textbf{\Large Mini-Project Report: Task 1}\\
\textbf{Preliminary Database Modeling for a Medical Database}
\end{center}

\vspace{0.8cm}

\textbf{Course Title:} Database Management System\\
\textbf{Course Code:} CSE 414

\vspace{0.8cm}

\textbf{Submitted to:}\\
Dr. Rudra Pratap Deb Nath\\
Professor\\
Department of Computer Science and Engineering\\
University of Chittagong

\vspace{0.8cm}

\textbf{Submitted by:}\\
Name: Sheikh Tawsif Bin Aftab\\
ID: 24701026\\
Session: 2023--2024 B.Sc. Engineering (4th Semester)\\
Department of Computer Science and Engineering\\
University of Chittagong

\vspace{0.8cm}

\textbf{Date of Submission:}\\
10 April 2026

\newpage

\section{Introduction}
The goal of this task is to prepare a preliminary database design for a medical information system.
The database must store information about drugs, their categories, diseases treated by those drugs,
side effects, commercial products, interactions with other drugs, and related clinical trials.

The design in this report was guided by the provided Excel sample. That workbook is not a valid
database table because it contains repeated groups such as multiple side-effect columns, multiple
interaction columns, and multiple clinical-trial-condition columns. Therefore, the first step was to
convert those repeating spreadsheet fields into proper entities and many-to-many relationships.

\section{What The Sample Data Shows}
After inspecting the sample workbook, the following observations were useful for the preliminary
design:

\begin{itemize}[leftmargin=1.5em]
    \item The workbook contains 250 rows.
    \item It contains 6 distinct named drugs and 49 diseases.
    \item Several spreadsheet columns are multi-valued: side effects, interacting drugs, and clinical trial conditions.
    \item A drug can belong to more than one category. For example, \texttt{abciximab} and \texttt{etanercept} appear under multiple drug categories.
    \item A drug can treat many diseases, and a disease can be treated by many drugs.
    \item A drug can have many commercial products, and those products may come from different companies.
    \item Clinical trial information is repeated across many rows, which indicates that trials should be stored separately and linked to drugs and conditions.
\end{itemize}

\section{Assumptions About The Mini-World}
\begin{enumerate}[leftmargin=1.5em]
    \item Each drug name in the spreadsheet represents a drug that should appear once in the \texttt{drugs} table.
    \item Interacting drugs are also stored as drugs, even if the sample file does not provide full details for them.
    \item Each disease belongs to one disease category in the current mini-world.
    \item Each product belongs to one drug and is marketed by one company.
    \item Each clinical trial has one main researcher and one main institution in this preliminary design.
    \item Clinical trial conditions are stored separately because the spreadsheet shows that one trial may mention multiple conditions.
\end{enumerate}

\section{Preliminary Database Description}
The preliminary design is centered on the \texttt{drugs} table. Around it, the model stores:

\begin{itemize}[leftmargin=1.5em]
    \item \textbf{Drug categories} through a bridge table, because a drug may belong to several categories.
    \item \textbf{Diseases} and \textbf{disease categories}, because drugs are used to treat diseases and many sample queries refer to disease categories.
    \item \textbf{Side effects}, because a drug may cause multiple side effects and the same side effect may occur for many drugs.
    \item \textbf{Products} and \textbf{companies}, because the same drug may be sold under several product names.
    \item \textbf{Drug interactions}, modeled as a self-relationship between drugs.
    \item \textbf{Clinical trials}, together with institutions, researchers, and trial conditions.
\end{itemize}

This design is already normalized enough to avoid the repeated columns seen in the workbook, but it
is still a preliminary Task 1 design. In Task 2, the model will be revised more formally using
course notation, explicit cardinalities, and stronger native-key choices.

\section{List Of Tables And Attributes}
\renewcommand{\arraystretch}{1.18}
\begin{longtable}{p{0.24\textwidth} p{0.70\textwidth}}
\toprule
\textbf{Table} & \textbf{Important attributes and purpose} \\
\midrule
\endfirsthead
\toprule
\textbf{Table} & \textbf{Important attributes and purpose} \\
\midrule
\endhead
\texttt{drugs} & \texttt{drug\_id} (PK), \texttt{drug\_name} (unique), \texttt{notes}. Stores each drug once. \\
\texttt{drug\_categories} & \texttt{drug\_category\_id} (PK), \texttt{category\_name} (unique). Stores available drug categories. \\
\texttt{drug\_category\_assignments} & \texttt{drug\_id} (FK), \texttt{drug\_category\_id} (FK). Resolves the many-to-many relation between drugs and drug categories. \\
\texttt{disease\_categories} & \texttt{disease\_category\_id} (PK), \texttt{category\_name} (unique). Stores categories such as Immunological or Respiratory. \\
\texttt{diseases} & \texttt{disease\_id} (PK), \texttt{disease\_name} (unique), \texttt{disease\_category\_id} (FK). Stores treatable diseases. \\
\texttt{side\_effects} & \texttt{side\_effect\_id} (PK), \texttt{side\_effect\_name} (unique). Stores side-effect names. \\
\texttt{drug\_treatments} & \texttt{drug\_id} (FK), \texttt{disease\_id} (FK). Stores which drugs treat which diseases. \\
\texttt{drug\_side\_effects} & \texttt{drug\_id} (FK), \texttt{side\_effect\_id} (FK). Stores side effects associated with each drug. \\
\texttt{companies} & \texttt{company\_id} (PK), \texttt{company\_name} (unique). Stores manufacturers or product owners. \\
\texttt{products} & \texttt{product\_id} (PK), \texttt{product\_name}, \texttt{drug\_id} (FK), \texttt{company\_id} (FK). Stores commercial products for drugs. \\
\texttt{drug\_interactions} & \texttt{drug\_id} (FK), \texttt{interacting\_drug\_id} (FK). Stores interacting pairs of drugs. \\
\texttt{institutions} & \texttt{institution\_id} (PK), \texttt{institution\_name}, \texttt{country}, \texttt{facility\_address}. Stores the trial host institution. \\
\texttt{researchers} & \texttt{researcher\_id} (PK), \texttt{researcher\_name}, \texttt{institution\_id} (FK). Stores main researchers. \\
\texttt{clinical\_trials} & \texttt{trial\_id} (PK), \texttt{trial\_title} (unique), \texttt{start\_date}, \texttt{completion\_date}, \texttt{participants}, \texttt{status}, \texttt{institution\_id} (FK), \texttt{main\_researcher\_id} (FK). Stores trial metadata. \\
\texttt{clinical\_trial\_drugs} & \texttt{trial\_id} (FK), \texttt{drug\_id} (FK). Stores which drugs are studied in a trial. \\
\texttt{trial\_conditions} & \texttt{trial\_condition\_id} (PK), \texttt{condition\_name} (unique). Stores conditions mentioned by trials. \\
\texttt{clinical\_trial\_conditions} & \texttt{trial\_id} (FK), \texttt{trial\_condition\_id} (FK). Stores the many-to-many relation between trials and conditions. \\
\bottomrule
\end{longtable}

\section{Why This Preliminary Design Can Answer The Sample Queries}
The design supports the sample queries because:

\begin{itemize}[leftmargin=1.5em]
    \item side-effect queries use \texttt{drug\_side\_effects} and \texttt{side\_effects},
    \item interaction queries use \texttt{drug\_interactions},
    \item disease-category queries use \texttt{drug\_treatments}, \texttt{diseases}, and \texttt{disease\_categories},
    \item product and company information come from \texttt{products} and \texttt{companies},
    \item clinical-trial queries use \texttt{clinical\_trials}, \texttt{clinical\_trial\_drugs}, \texttt{researchers}, and \texttt{institutions}.
\end{itemize}

\section{SQL Commands To Create The Database}
The complete SQL DDL for the preliminary design is included below. The same file is also stored in:

\texttt{SQL/Medical\_Database/task1\_preliminary\_schema.sql}

\lstinputlisting{../../SQL/Medical_Database/task1_preliminary_schema.sql}

\section{Conclusion}
This preliminary Task 1 database design converts the denormalized spreadsheet into a set of related
tables that can store the required medical information without repeated columns. The design keeps
the main concepts clear and supports the sample queries, while leaving room for improvement in the
more formal Task 2 redesign.

\end{document}
